% Part: counterfactuals
% Chapter: minimal-change-semantics
% Section: transitivity

\documentclass[../../../include/open-logic-section]{subfiles}

\begin{document}

\olfileid{cnt}{min}{cpo}

\olsection{Kontraposisi}

Kondisional material lan kondisional ketat ekuivalen karo
kontrapositife. Kontrafaktual ora mangkono. Iki tuladha saka
Kratzer:
\begin{quote}
  Saupama Goethe ora seda ing taun 1832, dheweke (tetep)
  mesthi wis seda saiki.

  Saupama Goethe durung seda saiki, dheweke mesthi wis seda
  ing taun 1832.
\end{quote}
Ukara kapisan bener: manungsa ora urip nganti atusan taun.
Ukara kapindho cetha luput: saupama Goethe durung seda saiki,
dheweke isih urip, mula ora mungkin wis seda ing taun 1832.

\begin{ex}\ollabel{ex:contraposition-counterex}
  Semantik bola nuduhake yen kontraposisi ora sah, yaiku $p
  \cif q \Entails/ \lnot q \cif \lnot p$. Upamane $p$
  ateges ``Goethe ora seda ing taun 1832'' lan $q$ ateges
  ``Goethe wis seda saiki.'' Iki bisa diwakili dening model
  $\mModel{M} = \tuple{W, O, V}$ kanthi $W = \{w, w_1, w_2\}$,
  $O_w = \{\{w\}, \{w, w_1\}, \{w, w_1, w_2\}\}$,
  $V(p) = \{w_1, w_2\}$ lan $V(q) = \{w, w_1\}$. Dadi $w$
  donya aktual kang ing kono Goethe seda ing taun 1832 lan
  saiki isih seda; $w_1$ donya kang cedhak, kang ing kono
  Goethe seda, upamane, ing taun 1833 lan saiki isih seda;
  lan $w_2$ donya kang adoh, kang ing kono Goethe isih urip.
\begin{figure}
\begin{center}
\begin{tikzpicture}[scale=.6,layerwidth=1.5]\tiny
  \spheresystem{3}
  \begin{scope}
    \clip \propositionplot{0}{.8}{50}{4.5} -- (4.5,-4.5) -- (-4.5,-4.5) -- (-4.5,4.5) -- (4.5,4.5);
    \spherelayer{3}
  \end{scope}
  \propositionintersect{0}{2}{110}{5.5}
  \proposition{0}{.8}{50}{4.5}
  \path (0,0) node[world] {$w$};
  \spherepos{0}{2}{node[world] {$w_1$}}
  \spherepos{-50}{3}{node[world] {$w_2$}}
  \spherepos{28}{3}{node {$q$}}
  \spherepos{42}{3}{node {$\lnot q$}}
  \spherepos{-55}{4}{node {$p$}}
  \spherepos{-67}{4}{node {$\lnot p$}}
\end{tikzpicture}
\caption{Tuladha panyanggah tumrap kontraposisi}
\ollabel{fig:contraposition}
\end{center}
\end{figure}
Ana bola kang ngakoni $p$, yaiku $S = \{w, w_1\}$, lan
$p \lif q$ bener ing kabeh donya ing kono, mula
$\mSat{M}{p \cif q}[w]$. Nanging bola siji-sijine kang
ngakoni $\lnot q$, yaiku $\{w, w_1, w_2\}$, ngemot donya~$w_2$
kang ing kono $q$ luput lan $p$ bener, mula
$\mSat/{M}{\lnot q \lif \lnot p}[w_2]$. Dadi
kontrafaktual kontrapositif luput ing~$w$.
\end{ex}

\end{document}
